In questo capitolo definiamo la parametrizzazione STERA delle proiezioni di un Transformer e ne ricaviamo le proprietà, restando sul piano della teoria. La \autoref{sec:proiezione-stera} dà la formula di una proiezione STERA, che combina quattro componenti: le coordinate dei blocchi RowMask23, una scala per blocco, i guadagni per riga e per colonna e un termine low-rank. La \autoref{sec:esempio-128} mostra come le quattro componenti costruiscono una proiezione $128\times128$. Le sezioni successive analizzano una componente alla volta: le coordinate e le proprietà del blocco RowMask23 (\autoref{sec:coordinate}), le scale (\autoref{sec:scale}), i guadagni (\autoref{sec:guadagni}) e il termine low-rank (\autoref{sec:basso-rango}). Il capitolo si chiude con il bilancio complessivo dei parametri. Come si calcolano il forward e la backpropagation di una proiezione STERA è descritto nel \autoref{cap:capitoli/04c-forward-backward}, il motore che li esegue nel \autoref{cap:capitoli/04d-motore}; i valori appresi dal modello addestrato e lo studio dell'importanza delle componenti nel \autoref{cap:capitoli/04-architettura}.

\section{La proiezione STERA}
\label{sec:proiezione-stera}

\begin{definizione}[Proiezione STERA]
\label{def:proiezione-stera}
Siano $m$ e $n$ multipli di 128 e $P=m/128$, $Q=n/128$. Una proiezione STERA $W\in\mathbb{R}^{m\times n}$ è
\[
W \;=\; \operatorname{diag}(g_{\mathrm{out}})\;\widetilde W\;\operatorname{diag}(g_{\mathrm{in}}) \;+\; U V^{\top},
\qquad
\widetilde W \;=\; \Big[\,\alpha^{(p,q)}\,\Phi\big(c^{(p,q)}\big)\Big]_{p=1,\dots,P;\;q=1,\dots,Q},
\]
dove $\widetilde W$ è la matrice a mosaico di $P\times Q$ blocchi $128\times128$ e $\Phi:\mathbb{R}^{3176}\to\mathbb{R}^{128\times128}$ è la mappa che costruisce un blocco RowMask23 dalle sue coordinate.
\end{definizione}

La matrice è quindi divisa in blocchi $128\times128$: ogni blocco $(p,q)$ è un blocco RowMask23 con coordinate $c^{(p,q)}$ proprie, moltiplicato per una scala $\alpha^{(p,q)}$. Sul mosaico così ottenuto agiscono due componenti che attraversano i confini dei blocchi: i guadagni $g_{\mathrm{out}}$ e $g_{\mathrm{in}}$, che moltiplicano ogni riga e ogni colonna dell'intera matrice, e il termine low-rank $UV^\top$, che si somma all'intera proiezione. La \autoref{tab:componenti-proiezione} riassume le quattro componenti, il numero di valori che ciascuna memorizza e la sezione che la tratta.

\begin{table}[H]
\centering
\small
\begin{tabular}{llrl}
\toprule
Componente & Simbolo & Valori & Sezione \\
\midrule
Coordinate & $c^{(p,q)}\in\mathbb{R}^{3176}$ & $3176\,PQ$ & \ref{sec:coordinate} \\
Scale & $\alpha^{(p,q)}\in\mathbb{R}$ & $PQ$ & \ref{sec:scale} \\
Guadagni & $g_{\mathrm{out}}\in\mathbb{R}^{m}$, $g_{\mathrm{in}}\in\mathbb{R}^{n}$ & $m+n$ & \ref{sec:guadagni} \\
Low-rank & $U\in\mathbb{R}^{m\times r}$, $V\in\mathbb{R}^{n\times r}$ & $r(m+n)$ & \ref{sec:basso-rango} \\
\midrule
Bilancio complessivo & & & \ref{sec:bilancio-parametri} \\
\bottomrule
\end{tabular}
\caption{Le quattro componenti di una proiezione STERA $m\times n$ con termine low-rank di rango $r$; $PQ=mn/16384$ è il numero di blocchi.}
\label{tab:componenti-proiezione}
\end{table}

\section{Una proiezione $128\times128$}
\label{sec:esempio-128}

Il caso più semplice è una proiezione formata da un solo blocco, $m=n=128$ e $P=Q=1$. La formula della \autoref{def:proiezione-stera} diventa
\[
W \;=\; \operatorname{diag}(g_{\mathrm{out}})\;\alpha\,\Phi(c)\;\operatorname{diag}(g_{\mathrm{in}}) \;+\; UV^\top,
\]
e un singolo elemento della matrice vale
\begin{equation}
W_{ij} \;=\; g_{\mathrm{out},i}\;\cdot\;\alpha\;\cdot\;\Phi(c)_{ij}\;\cdot\;g_{\mathrm{in},j}\;+\;\sum_{k=1}^{r} U_{ik}V_{jk}.
\label{eq:elemento-proiezione}
\end{equation}
Ogni componente interviene quindi in un modo diverso sullo stesso elemento: le coordinate ne fissano il valore di partenza, la scala lo moltiplica insieme a tutto il blocco, i guadagni lo moltiplicano insieme alla sua riga e alla sua colonna, il low-rank gli somma un contributo. La \autoref{tab:proiezione-128}, subito sotto, lo riassume.

\begin{table}[H]
\centering
\small
\renewcommand{\arraystretch}{1.25}
\begin{tabular}{>{\raggedright\arraybackslash}p{24mm}>{\raggedright\arraybackslash}p{62mm}r}
\toprule
Componente & Effetto sull'elemento $(i,j)$ & Valori \\
\midrule
\rowcolor{steraBlu!12} Coordinate $c$ & valore di partenza $\Phi(c)_{ij}$: $d_i$, $b_{i\oplus j}$ oppure $b_{i\oplus j}+e_{i,k}$ & \num{3176} \\
\rowcolor{steraGrigio!12} Scala $\alpha$ & moltiplica tutti gli elementi del blocco & $1$ \\
\rowcolor{steraVerde!15} Guadagni & moltiplicano per $g_{\mathrm{out},i}$ (riga) e $g_{\mathrm{in},j}$ (colonna) & $256$ \\
\rowcolor{steraArancio!15} Low-rank & somma $\sum_{k} U_{ik}V_{jk}$ & $256\,r$ \\
\bottomrule
\end{tabular}
\caption{Come le quattro componenti di una proiezione STERA $128\times128$ determinano l'elemento $(i,j)$, secondo la \eqref{eq:elemento-proiezione}, e quanti valori memorizza ciascuna.}
\label{tab:proiezione-128}
\end{table}

L'esempio mostra anche perché il low-rank va dimensionato rispetto alla matrice. In una proiezione $128\times128$ coordinate, scala e guadagni costano \num{3433} valori contro i \num{16384} della forma densa, mentre il low-rank ne aggiunge $256$ per ogni unità di rango: la proiezione resta più piccola della forma densa solo se $r<50$. Nelle proiezioni di STERA 30M, che hanno almeno 768 righe e colonne, il rango 96 costa invece una frazione molto minore, come mostra il bilancio della \autoref{sec:bilancio-parametri}.

% ---------------------------------------------------------------- COORDINATE

\section{Le coordinate: il blocco RowMask23}
\label{sec:coordinate}

Le coordinate sono il contributo principale della tesi. Questa sezione introduce lo XOR fra indici, definisce il blocco RowMask23 e ne ricava le proprietà.

\subsection{Lo XOR fra indici e le matrici XOR-circolanti}
\label{sec:notazione}

\paragraph{Notazione.} Fissiamo $N=2^n$ con $n=7$, cioè $N=128$. Gli indici di riga e di colonna vanno da $0$ a $N-1$ e li identifichiamo con la loro rappresentazione binaria a $n$ bit. Indichiamo con $\oplus$ lo XOR bit a bit:
\[
i\oplus j = \sum_{\ell=0}^{n-1} \big((i_\ell + j_\ell)\bmod 2\big)\,2^\ell,
\]
dove $i_\ell$ e $j_\ell$ sono i bit di posizione $\ell$. L'insieme $\{0,\dots,N-1\}$ con l'operazione $\oplus$ è un gruppo abeliano, isomorfo a $(\mathbb{Z}_2)^n$, con elemento neutro $0$; ogni elemento è l'inverso di sé stesso ($i\oplus i=0$). Ne seguono due fatti che useremo spesso:
\begin{itemize}
\item per $i$ fissato, la funzione $j\mapsto i\oplus j$ è una biiezione di $\{0,\dots,N-1\}$ in sé;
\item $i\oplus j = s$ se e solo se $j = i\oplus s$.
\end{itemize}

\begin{definizione}[Matrice XOR-circolante]
Dato un vettore $b\in\mathbb{R}^N$, la matrice XOR-circolante generata da $b$ è la matrice $B\in\mathbb{R}^{N\times N}$ con elementi $B_{ij} = b_{i\oplus j}$.
\end{definizione}

Il nome richiama le matrici circolanti, in cui l'elemento $(i,j)$ dipende da $(j-i)\bmod N$; qui dipende invece da $i\oplus j$. In letteratura queste matrici sono dette anche \emph{diadiche} \cite{beauchamp1984applications}. Ogni riga e ogni colonna di $B$ contiene tutti gli elementi di $b$, ciascuno una sola volta, in un ordine che cambia da riga a riga. La matrice è simmetrica, perché $i\oplus j=j\oplus i$, e $b_0$ compare esattamente sulla diagonale principale.

$B$ si può vedere come l'unione di $N$ \emph{diagonali generalizzate}: la diagonale di indice $s$ è l'insieme delle posizioni $\{(i,i\oplus s): i=0,\dots,N-1\}$, su cui la matrice è costante e vale $b_s$. Le $N$ diagonali generalizzate formano una partizione delle $N^2$ posizioni. Una matrice XOR-circolante ha però solo $N$ gradi di libertà su $N^2$ posizioni, troppo pochi per una rete che deve imparare trasformazioni generali: RowMask23 parte da questa base e le aggiunge una diagonale libera e un insieme fisso di correzioni per riga.

\subsection{Il blocco RowMask23}
\label{sec:blocco-rowmask}

\subsubsection{Supporto delle correzioni}

Per ogni riga $i\in\{1,\dots,N-1\}$ definiamo l'insieme di 23 offset
\begin{equation}
S_i = \big\{\,\sigma(i,k) : k=0,\dots,22\,\big\},\qquad
\sigma(i,k) = 1 + \big((13k + 7i)\bmod 127\big).
\label{eq:supporto}
\end{equation}
Ogni offset è compreso tra $1$ e $127$, quindi diverso da zero. La riga $0$ non ha correzioni: $S_0=\emptyset$.

\begin{proposizione}
Per ogni riga $i\ge1$, i 23 offset $\sigma(i,0),\dots,\sigma(i,22)$ sono distinti.
\end{proposizione}

\begin{proof}
Se $\sigma(i,k)=\sigma(i,k')$ con $0\le k,k'\le22$, allora $13(k-k')\equiv0\pmod{127}$. Poiché $127$ è primo e non divide $13$, deve dividere $k-k'$; ma $|k-k'|\le22<127$, quindi $k=k'$.
\end{proof}

I coefficienti $13$ e $7$ distribuiscono le correzioni in modo che righe diverse usino offset diversi, evitando che alcune diagonali generalizzate ricevano molte correzioni e altre nessuna: lo spostamento $7i$ fa ruotare l'insieme da una riga all'altra, il passo $13$ distanzia gli offset dentro la stessa riga.

\subsubsection{Definizione}

\begin{definizione}[Blocco RowMask23]
Il vettore di coordinate di un blocco è $c=(d,b,e)$ con
\begin{itemize}
\item $d\in\mathbb{R}^{128}$, la \emph{diagonale}, con componenti $d_i$, $i=0,\dots,127$;
\item $b\in\mathbb{R}^{127}$, la \emph{base XOR}, con componenti $b_s$, $s=1,\dots,127$;
\item $e\in\mathbb{R}^{127\times23}$, le \emph{correzioni}, con componenti $e_{i,k}$, $i=1,\dots,127$, $k=0,\dots,22$.
\end{itemize}
Il blocco $W=\Phi(c)\in\mathbb{R}^{128\times128}$ ha elementi
\begin{equation}
W_{ij} =
\begin{cases}
d_i & \text{se } i=j,\\[2pt]
b_{i\oplus j} + e_{i,k} & \text{se } i\ne j \text{ e } i\oplus j=\sigma(i,k) \text{ per qualche } k,\\[2pt]
b_{i\oplus j} & \text{altrimenti.}
\end{cases}
\label{eq:rowmask23}
\end{equation}
\end{definizione}

In parole: la diagonale è libera; fuori dalla diagonale ogni elemento eredita il valore della base XOR della sua diagonale generalizzata; in 23 posizioni per riga, scelte da \eqref{eq:supporto}, a questo valore si aggiunge una correzione propria. Il nome ``RowMask23'' indica questa maschera di 23 posizioni per riga. Equivalentemente,
\begin{equation}
W = D + B + R,
\label{eq:scomposizione}
\end{equation}
dove $D=\mathrm{diag}(d)$, $B$ è la matrice XOR-circolante generata da $(0,b_1,\dots,b_{127})$ e $R$ è una matrice sparsa con \num{2921} elementi non nulli, $R_{i,\,i\oplus\sigma(i,k)}=e_{i,k}$.

\paragraph{Lettura del blocco.} Il blocco si legge in tre passi, ciascuno dei quali occupa posizioni diverse della matrice; la \autoref{tab:struttura-blocco}, subito sotto, mostra le prime otto righe e colonne di un blocco con un colore per ogni passo.
\begin{enumerate}
\item \textbf{La diagonale} (in blu). I 128 elementi $(i,i)$ sono liberi: ciascuno ha un proprio valore $d_i$.
\item \textbf{La base XOR} (in grigio). Ogni elemento fuori diagonale $(i,j)$ prende il valore $b_{i\oplus j}$, dove $\oplus$ è lo XOR bit a bit degli indici. Tutte le posizioni con lo stesso valore di $i\oplus j$ formano una \emph{diagonale generalizzata} e condividono lo stesso valore: nella tabella sono le caselle con lo stesso indice, per esempio tutte quelle con $b_1$. Bastano così 127 valori $b_1,\dots,b_{127}$ per riempire i \num{16256} elementi fuori diagonale.
\item \textbf{Le correzioni} (in arancione). In 23 posizioni fisse di ogni riga, dalla riga 1 alla 127, al valore della base si somma una correzione propria $e_{i,k}$, che rende quella posizione diversa dalle altre della sua diagonale generalizzata: \num{2921} valori in tutto.
\end{enumerate}
Le $128+127+\num{2921}=\num{3176}$ \emph{coordinate} $c=(d,b,e)$ determinano l'intero blocco, che indichiamo con $\Phi(c)$, al posto dei \num{16384} valori di una matrice densa.


\begin{table}[H]
\centering
\footnotesize
\setlength{\tabcolsep}{3pt}
\renewcommand{\arraystretch}{1.35}
\begin{tabular}{c|*{8}{>{\centering\arraybackslash}p{9.5mm}}}
$i\backslash j$ & \textbf{0} & \textbf{1} & \textbf{2} & \textbf{3} & \textbf{4} & \textbf{5} & \textbf{6} & \textbf{7} \\
\hline
\textbf{0} & \cellcolor{steraBlu!30}$d_{0}$ & \cellcolor{steraGrigio!12}$b_{1}$ & \cellcolor{steraGrigio!12}$b_{2}$ & \cellcolor{steraGrigio!12}$b_{3}$ & \cellcolor{steraGrigio!12}$b_{4}$ & \cellcolor{steraGrigio!12}$b_{5}$ & \cellcolor{steraGrigio!12}$b_{6}$ & \cellcolor{steraGrigio!12}$b_{7}$ \\
\textbf{1} & \cellcolor{steraArancio!40}$b_{1}{+}e$ & \cellcolor{steraBlu!30}$d_{1}$ & \cellcolor{steraGrigio!12}$b_{3}$ & \cellcolor{steraGrigio!12}$b_{2}$ & \cellcolor{steraGrigio!12}$b_{5}$ & \cellcolor{steraGrigio!12}$b_{4}$ & \cellcolor{steraGrigio!12}$b_{7}$ & \cellcolor{steraGrigio!12}$b_{6}$ \\
\textbf{2} & \cellcolor{steraGrigio!12}$b_{2}$ & \cellcolor{steraGrigio!12}$b_{3}$ & \cellcolor{steraBlu!30}$d_{2}$ & \cellcolor{steraGrigio!12}$b_{1}$ & \cellcolor{steraGrigio!12}$b_{6}$ & \cellcolor{steraGrigio!12}$b_{7}$ & \cellcolor{steraGrigio!12}$b_{4}$ & \cellcolor{steraArancio!40}$b_{5}{+}e$ \\
\textbf{3} & \cellcolor{steraGrigio!12}$b_{3}$ & \cellcolor{steraArancio!40}$b_{2}{+}e$ & \cellcolor{steraGrigio!12}$b_{1}$ & \cellcolor{steraBlu!30}$d_{3}$ & \cellcolor{steraGrigio!12}$b_{7}$ & \cellcolor{steraGrigio!12}$b_{6}$ & \cellcolor{steraGrigio!12}$b_{5}$ & \cellcolor{steraGrigio!12}$b_{4}$ \\
\textbf{4} & \cellcolor{steraGrigio!12}$b_{4}$ & \cellcolor{steraGrigio!12}$b_{5}$ & \cellcolor{steraArancio!40}$b_{6}{+}e$ & \cellcolor{steraGrigio!12}$b_{7}$ & \cellcolor{steraBlu!30}$d_{4}$ & \cellcolor{steraGrigio!12}$b_{1}$ & \cellcolor{steraGrigio!12}$b_{2}$ & \cellcolor{steraGrigio!12}$b_{3}$ \\
\textbf{5} & \cellcolor{steraGrigio!12}$b_{5}$ & \cellcolor{steraGrigio!12}$b_{4}$ & \cellcolor{steraGrigio!12}$b_{7}$ & \cellcolor{steraGrigio!12}$b_{6}$ & \cellcolor{steraGrigio!12}$b_{1}$ & \cellcolor{steraBlu!30}$d_{5}$ & \cellcolor{steraArancio!40}$b_{3}{+}e$ & \cellcolor{steraGrigio!12}$b_{2}$ \\
\textbf{6} & \cellcolor{steraGrigio!12}$b_{6}$ & \cellcolor{steraArancio!40}$b_{7}{+}e$ & \cellcolor{steraGrigio!12}$b_{4}$ & \cellcolor{steraGrigio!12}$b_{5}$ & \cellcolor{steraGrigio!12}$b_{2}$ & \cellcolor{steraGrigio!12}$b_{3}$ & \cellcolor{steraBlu!30}$d_{6}$ & \cellcolor{steraGrigio!12}$b_{1}$ \\
\textbf{7} & \cellcolor{steraGrigio!12}$b_{7}$ & \cellcolor{steraGrigio!12}$b_{6}$ & \cellcolor{steraGrigio!12}$b_{5}$ & \cellcolor{steraArancio!40}$b_{4}{+}e$ & \cellcolor{steraGrigio!12}$b_{3}$ & \cellcolor{steraGrigio!12}$b_{2}$ & \cellcolor{steraArancio!40}$b_{1}{+}e$ & \cellcolor{steraBlu!30}$d_{7}$ \\
\end{tabular}

\medskip
\begin{tabular}{lll}
\cellcolor{steraBlu!30}\phantom{xx} diagonale $d_i$ &
\cellcolor{steraGrigio!12}\phantom{xx} base XOR $b_{i\oplus j}$ &
\cellcolor{steraArancio!40}\phantom{xx} base più correzione $b_{i\oplus j}+e_{i,k}$
\end{tabular}
\caption{Le prime 8 righe e 8 colonne di un blocco RowMask23 $128\times128$. Ogni casella indica da quali coordinate dipende: la diagonale è libera, le altre caselle prendono la coordinata di base di indice $i\oplus j$ e, nelle posizioni scelte dalla maschera, anche una correzione propria. Caselle con lo stesso indice, per esempio tutte quelle con $b_1$, appartengono alla stessa diagonale generalizzata e condividono il valore della base. La riga 0 non ha correzioni.}
\label{tab:struttura-blocco}
\end{table}

\subsubsection{Un esempio}
\label{sec:guida-visiva-rowmask}
\label{sec:esempio-numerico-rowmask}

Seguiamo una sola coordinata di base, $b_8$, in due righe diverse del blocco.

Si osservi la coordinata $b_8$. Nella posizione $(0,8)$ vale $0\oplus8=8$, quindi l'elemento riceve $b_8$; anche nella posizione $(1,9)$ vale $1\oplus9=8$. La riga $1$, però, per $k=0$ ha offset $\sigma(1,0)=1+(7\bmod127)=8$, quindi $(1,9)$ è una posizione con correzione, mentre la riga $0$ non ne ha. Assegnando a scopo illustrativo i valori $d_1=0{,}60$, $b_8=0{,}25$ ed $e_{1,0}=-0{,}10$, che non sono pesi del modello addestrato, si ottiene la \autoref{fig:lettura-casella-rowmask}.

\begin{table}[H]
\centering
\small
\begin{tabular}{cccl}
\toprule
Casella $(i,j)$ & $s=i\oplus j$ & Origine & Peso risultante \\
\midrule
$(1,1)$ & $0$ & diagonale & $d_1 = 0{,}60$ \\
$(0,8)$ & $8$ & base, riga senza correzioni & $b_8 = 0{,}25$ \\
$(1,9)$ & $8$ & base e correzione $k=0$ & $b_8 + e_{1,0} = 0{,}15$ \\
\bottomrule
\end{tabular}
\caption{Lettura di tre caselle di un blocco RowMask23, con valori dimostrativi. Le caselle $(0,8)$ e $(1,9)$ condividono la coordinata di base $b_8$; solo la seconda riceve anche una correzione.}
\label{fig:lettura-casella-rowmask}
\end{table}

L'esempio mostra l'idea del blocco: la coordinata $b_8$ è \emph{condivisa} da tutta la sua diagonale generalizzata, e la correzione $e_{1,0}$ permette di \emph{adattare localmente} una sola casella senza toccare le altre. Il modello non memorizza i tre elementi come pesi indipendenti, ma le coordinate $d_1$, $b_8$ ed $e_{1,0}$.

\label{sec:non-potatura-rowmask}
La parola \emph{maschera} non indica una matrice quasi tutta nulla. Quasi tutti gli elementi fuori diagonale dipendono da una coordinata della base e sono in generale diversi da zero; la parte sparsa riguarda solo le correzioni. RowMask23 non è quindi una potatura della matrice: è una matrice piena con pochi gradi di libertà.

\subsection{Proprietà del blocco}
\label{sec:linearita-iniettivita}

\subsubsection{Gradi di libertà}

\begin{proposizione}
Un blocco RowMask23 ha
\[
C = 128 + 127 + 127\cdot23 = \num{3176}
\]
coordinate, contro i \num{16384} elementi di una matrice densa $128\times128$. Il rapporto è $\num{3176}/\num{16384}\approx 0{,}1938$, cioè l'$80{,}6\%$ di gradi di libertà in meno, un fattore $5{,}16$.
\end{proposizione}

Il rapporto non dipende dalla precisione numerica, purché coordinate e matrice densa usino la stessa. Le correzioni sono la parte più numerosa, il $92\%$ delle coordinate; la base XOR occupa meno del $4\%$ del blocco ma tocca tutti i \num{16256} elementi fuori diagonale.

\subsubsection{Linearità e iniettività}

\begin{proposizione}
La mappa $\Phi:\mathbb{R}^{3176}\to\mathbb{R}^{128\times128}$ è lineare e iniettiva. La sua immagine, la \emph{famiglia RowMask23}, è quindi un sottospazio lineare di dimensione \num{3176} dello spazio delle matrici $128\times128$.
\end{proposizione}

\begin{proof}
La linearità segue da \eqref{eq:rowmask23}: ogni elemento di $W$ è la somma di al più due coordinate con coefficiente uno. Per l'iniettività basta leggere tutte le coordinate dagli elementi di $W$. La diagonale dà $d_i=W_{ii}$. La riga $0$ non ha correzioni, quindi $W_{0,s}=b_s$ per $s\ge1$ e la prima riga fornisce l'intera base. Infine, per $i\ge1$, la posizione $j=i\oplus\sigma(i,k)$ dà $e_{i,k}=W_{ij}-b_{\sigma(i,k)}$. Due vettori di coordinate diversi producono quindi matrici diverse.
\end{proof}

\subsubsection{Ricostruzione esatta}
\label{sec:esattezza-ricostruzione}

La proposizione precedente dà un significato preciso alla parola \emph{Reconstruction} nel nome STERA. Una matrice $128\times128$ con la struttura RowMask23, cioè che appartiene all'immagine di $\Phi$, si codifica con le sue \num{3176} coordinate e si ricostruisce da esse in modo esatto: la corrispondenza fra le coordinate e le matrici della famiglia è biunivoca e la ricostruzione non richiede di stimare coefficienti mancanti. La proprietà riguarda le matrici della famiglia: una matrice densa qualunque non ammette in generale una rappresentazione RowMask23, perché la famiglia ha dimensione \num{3176}, inferiore a \num{16384}. La compressione non è quindi l'approssimazione di una matrice densa, ma la scelta di lavorare dentro un sottospazio.

Questa corrispondenza è ciò che serve al training. L'optimizer non modifica una matrice densa da proiettare poi sulla famiglia: modifica le coordinate, e a ogni vettore di coordinate corrisponde una e una sola matrice. Ne seguono tre conseguenze:
\begin{itemize}
\item il training può raggiungere ogni matrice della famiglia, quindi conserva per intero l'espressività che la struttura RowMask23 mette a disposizione;
\item non esistono coordinate ridondanti, cioè direzioni di aggiornamento che lascino invariata la matrice: ogni coordinata contribuisce al modello;
\item il gradiente rispetto alle coordinate (\autoref{sec:gradiente-coordinate}) passa per la stessa $\Phi$ usata nel forward, quindi l'ottimizzazione migliora esattamente la matrice $\Phi(c)$ che il modello usa.
\end{itemize}
La stessa esattezza vale fra training e inferenza: la ricostruzione eseguita dai kernel di inferenza coincide bit per bit con quella del training (\autoref{sec:kernel-rowmask}), e il modello che si esegue è lo stesso che si è addestrato.

\subsubsection{Rango}

Una matrice con pochi gradi di libertà potrebbe sembrare degenere, come una fattorizzazione $UV^\top$ di rango basso. Non è così.

\begin{proposizione}
L'insieme dei vettori di coordinate $c$ per cui $\Phi(c)$ è singolare ha misura di Lebesgue nulla in $\mathbb{R}^{3176}$. In particolare, quasi ogni matrice RowMask23 ha rango pieno $128$.
\end{proposizione}

\begin{proof}
La funzione $p(c)=\det\Phi(c)$ è un polinomio nelle coordinate, perché il determinante è un polinomio negli elementi e questi sono lineari in $c$. Il polinomio non è identicamente nullo: con $d_i=1$ per ogni $i$ e tutte le altre coordinate nulle si ottiene l'identità. L'insieme degli zeri di un polinomio non nullo ha misura nulla.
\end{proof}

Numero di parametri e rango sono quindi cose distinte. Una fattorizzazione di rango 16 in dimensione 128 ha $2\cdot128\cdot16=\num{4096}$ parametri, più di RowMask23, ma ignora $112$ direzioni dell'ingresso; un blocco RowMask23 ne usa \num{3176} e in generale agisce su tutte le direzioni. Le due famiglie sono sottospazi diversi, con geometrie diverse: per questo nella proiezione STERA si possono combinare (\autoref{sec:basso-rango}). Allo stesso modo, rango pieno non significa coefficienti arbitrari: la famiglia RowMask23 può avere rango pieno e restare un sottoinsieme proprio delle matrici dense.


% ---------------------------------------------------------------- ALTRE COMPONENTI

\section{Le scale dei blocchi}
\label{sec:scale}

\subsection{Che cosa sono e con quale criterio si applicano}
Ogni blocco $(p,q)$ del mosaico è moltiplicato per un numero $\alpha^{(p,q)}$, la sua scala. Il criterio è uniforme: una scala per ogni blocco $128\times128$ di ogni proiezione, cioè una ogni \num{16384} valori della forma densa, per \num{4950} scale in STERA 30M. Le coordinate fissano la forma del blocco, la scala ne fissa l'ampiezza: separare le due cose permette a blocchi diversi della stessa proiezione di avere ampiezze molto diverse senza che l'optimizer debba spostare in modo coordinato tutte le \num{3176} coordinate.

\subsection{Inizializzazione e costo}
Le scale partono da $1$ e non sono soggette al weight decay: una penalità che le spingesse verso zero ridurrebbe proprio l'ampiezza che devono controllare. Costano $mn/16384$ valori per proiezione e una moltiplicazione per blocco nel prodotto.

\section{I guadagni per riga e per colonna}
\label{sec:guadagni}

\subsection{Che cosa sono e con quale criterio si applicano}
I guadagni sono due vettori, $g_{\mathrm{out}}\in\mathbb{R}^m$ e $g_{\mathrm{in}}\in\mathbb{R}^n$, che moltiplicano ogni riga e ogni colonna dell'intera matrice: $\big(\operatorname{diag}(g_{\mathrm{out}})\,\widetilde W\,\operatorname{diag}(g_{\mathrm{in}})\big)_{ij} = g_{\mathrm{out},i}\,\widetilde W_{ij}\,g_{\mathrm{in},j}$. Il criterio è uno per ogni riga e uno per ogni colonna di ogni proiezione. Le scale agiscono su un blocco alla volta, i guadagni attraversano i confini dei blocchi: danno a ogni componente in uscita e in ingresso un'ampiezza propria, uguale in tutti i blocchi che la contengono.

\subsection{Inizializzazione e costo}
I guadagni partono da $1$ e, come le scale, non sono soggetti al weight decay. Costano $m+n$ valori per proiezione; nel prodotto con un vettore equivalgono a moltiplicare l'ingresso per $g_{\mathrm{in}}$ e l'uscita per $g_{\mathrm{out}}$, con $m+n$ moltiplicazioni.

\section{Il termine low-rank}
\label{sec:basso-rango}

\subsection{Che cosa è e con quale criterio si applica}
Il termine $UV^\top$, con $U\in\mathbb{R}^{m\times r}$ e $V\in\mathbb{R}^{n\times r}$, è una matrice $m\times n$ di rango al più $r$ che si somma all'intera proiezione. Il criterio è lo stesso per tutte le proiezioni del modello, embedding compreso: in STERA 30M il rango è $r=96$, fissato prima del training. Il mosaico di blocchi RowMask23 tratta ogni blocco in modo indipendente, e ciascuno è vincolato alla struttura XOR con 23 correzioni per riga; il termine low-rank non ha né l'uno né l'altro vincolo, attraversa tutti i blocchi e può rappresentare un numero limitato di direzioni qualunque. Le due rappresentazioni sono complementari: RowMask23 dà una matrice a rango pieno con pochi valori, il low-rank dà poche direzioni senza struttura.

\subsection{Inizializzazione e costo}
$U$ parte da zero e $V$ da valori normali di deviazione standard $0{,}02$: all'inizio il termine è nullo e ogni proiezione coincide con il suo mosaico di blocchi, mentre il gradiente rispetto a $U$ è già diverso da zero. A differenza di scale e guadagni, $U$ e $V$ sono soggetti al weight decay, come le coordinate. Il termine memorizza $r(m+n)$ valori e il suo prodotto con un vettore si calcola come $U(V^\top x)$, con $r(m+n)$ moltiplicazioni, senza formare la matrice $m\times n$. In STERA 30M i fattori sono memorizzati e addestrati in FP32 come ogni altro parametro.

% ---------------------------------------------------------------- BILANCIO

\section{Bilancio dei parametri}
\label{sec:bilancio-parametri}

Il risparmio garantito dai blocchi RowMask23 riguarda la parte strutturata delle matrici. Contiamo ora tutte le componenti, prima per una proiezione, poi per un layer del Transformer e infine per un modello intero.

\begin{proposizione}[Valori memorizzati e risparmio di una proiezione]
\label{prop:risparmio-proiezione}
La proiezione della \autoref{def:proiezione-stera} memorizza
\[
P(m,n,r) \;=\; \frac{3177}{16384}\,mn \;+\; (1+r)\,(m+n)
\]
valori. Il risparmio rispetto ai $mn$ valori della forma densa è
\[
s(m,n,r) \;=\; 1-\frac{3177}{16384}-(1+r)\left(\frac1m+\frac1n\right)
\;\approx\; 80{,}61\% \;-\; (1+r)\left(\frac1m+\frac1n\right).
\]
\end{proposizione}

\begin{proof}
I blocchi sono $mn/16384$; ciascuno contribuisce con 3176 coordinate e una scala, cioè $3177\,mn/16384$ valori in tutto. I guadagni aggiungono $m+n$ valori e i fattori $U$ e $V$ ne aggiungono $r(m+n)$. Dividendo per $mn$ si ottiene $s$.
\end{proof}

La formula separa ciò che il formato garantisce da ciò che dipende dalle scelte di progetto:
\begin{itemize}
\item \textbf{coordinate e scale} occupano sempre la frazione $3177/16384=19{,}39\%$ della forma densa, qualunque siano le dimensioni: è il risparmio dell'$80{,}61\%$ garantito da RowMask23;
\item \textbf{i guadagni} dipendono solo dalla forma della matrice e pesano $1/m+1/n$, al più l'$1{,}56\%$ per una matrice $128\times128$;
\item \textbf{il low-rank} dipende dalla scelta di $r$ ed è l'unica componente che può spostare sensibilmente il risparmio.
\end{itemize}

\begin{osservazione}[Limiti]
Per ogni scelta dei parametri $s(m,n,r)<80{,}61\%$. In assenza del termine low-rank ($r=0$), e poiché $m,n\ge128$, vale $s(m,n,0)\ge 80{,}61\%-\tfrac{2}{128}=79{,}05\%$.
\end{osservazione}

\paragraph{Un layer del Transformer.} In un layer con dimensione $d$ e feed-forward $f=4d$, le proiezioni sono quattro matrici $d\times d$ (attention) e due matrici $4d\times d$ e $d\times 4d$ (feed-forward), per $12d^2$ valori densi. Se il rango è proporzionale alla dimensione, $r=\kappa d$, sommando la \autoref{prop:risparmio-proiezione} sulle sei matrici si ottiene
\[
s_{\text{layer}} \;=\; 1-\frac{3177}{16384} - 1{,}5\,\kappa - \frac{1{,}5}{d}
\;\approx\; 80{,}61\% - 1{,}5\,\kappa .
\]
Il coefficiente $1{,}5$ è la media, pesata per le dimensioni, del costo del low-rank nelle due famiglie di matrici: $2\kappa$ per ognuna delle quattro matrici $d\times d$ e $1{,}25\,\kappa$ per ognuna delle due matrici del feed-forward, cioè $(4\cdot 2\kappa\cdot d^2+2\cdot1{,}25\,\kappa\cdot4d^2)/12d^2=1{,}5\,\kappa$. Il termine $1{,}5/d$, dovuto ai guadagni, vale lo $0{,}2\%$ per $d=768$ e tende a zero. Ne seguono due comportamenti al crescere del modello:
\begin{itemize}
\item \textbf{con il rango proporzionale alla dimensione} ($\kappa$ fissato) il risparmio per layer è costante: con $\kappa=1/8$, come in STERA 30M ($96/768$), vale $61{,}9\%$ per ogni $d$;
\item \textbf{con il rango fissato} il termine del low-rank decresce come $1/d$ e il risparmio tende all'$80{,}61\%$ del formato.
\end{itemize}

\paragraph{Il modello intero.} Il modello aggiunge l'embedding condiviso con la output head, una proiezione $V\times d$ molto alta che risparmia più dei layer, e alcuni parametri non compressi, identici nelle due forme.

\begin{proposizione}[Risparmio di un modello STERA]
\label{prop:risparmio-modello}
Si consideri un Transformer decoder-only con $L$ layer, dimensione $d$, feed-forward $f$, vocabolario $V$ arrotondato a un multiplo di 128, tabella di $T$ posizioni ed embedding condiviso con la output head, in cui ogni proiezione è una proiezione STERA di rango $r$. Siano
\[
\begin{aligned}
M &= V d + L\,(4d^2 + 2fd) &&\text{(valori delle proiezioni nella forma densa)},\\
R &= (V+d) + L\,\big(8d + 2(f+d)\big) &&\text{(somma di righe e colonne delle proiezioni)},\\
A &= T d + 2d\,(2L+1) + L\,(2d+f) &&\text{(posizioni, LayerNorm e bias, non compressi)}.
\end{aligned}
\]
La forma densa equivalente conta $D = M + A$ valori, il modello STERA ne memorizza
\[
P \;=\; \underbrace{\frac{3177}{16384}\,M}_{\text{coordinate e scale}} \;+\; \underbrace{R}_{\text{guadagni}} \;+\; \underbrace{r\,R}_{\text{low-rank}} \;+\; \underbrace{A}_{\text{non compressi}},
\]
e il risparmio è
\[
S \;=\; 1 - \frac{P}{D} \;=\; \frac{M}{M+A}\left(1-\frac{3177}{16384}\right) \;-\; \frac{(1+r)\,R}{M+A}.
\]
\end{proposizione}

\begin{proof}
Si somma la \autoref{prop:risparmio-proiezione} sulle proiezioni del modello: l'embedding ($V\times d$) e, in ogni layer, quattro matrici $d\times d$ e due matrici $f\times d$ e $d\times f$. La somma dei prodotti $m_i n_i$ è $M$ e la somma di $m_i+n_i$ è $R$. Ai valori delle proiezioni si aggiungono, in entrambe le forme, gli $A$ parametri non compressi.
\end{proof}

La seconda forma di $S$ rende esplicito il ruolo di ogni componente: il primo termine è l'$80{,}61\%$ garantito da RowMask23, ridotto della quota $A/(M+A)$ di parametri che non sono matrici; il secondo è il costo di guadagni e low-rank, l'unico che dipende dal rango $r$. Per STERA 30M ($L=6$, $d=768$, $f=3072$, $V=\num{50304}$, $T=\num{2048}$, $r=96$) si ha $M=\num{81100800}$, $R=\num{134016}$, $A=\num{1620480}$ e quindi
\[
P = \num{15726150} + \num{134016} + \num{12865536} + \num{1620480} = \num{30346182},
\qquad D = \num{82721280},
\qquad S = 63{,}32\%.
\]
La \autoref{tab:componenti-stera30m} mostra lo stesso conto per componente: coordinate e low-rank insieme fanno oltre il $94\%$ dei valori memorizzati.

\begin{table}[H]
\centering
\small
\begin{tabular}{lrr}
\toprule
Componente di STERA 30M & Valori & Quota \\
\midrule
Coordinate RowMask23 (\num{4950} blocchi) & \num{15721200} & $51{,}8\%$ \\
Scale dei blocchi & \num{4950} & $0{,}02\%$ \\
Guadagni per riga e colonna & \num{134016} & $0{,}4\%$ \\
Termine low-rank ($r=96$) & \num{12865536} & $42{,}4\%$ \\
Posizioni, LayerNorm, bias & \num{1620480} & $5{,}3\%$ \\
\midrule
Totale & \num{30346182} & \\
Forma densa equivalente & \num{82721280} & \\
Risparmio & $63{,}32\%$ & \\
\bottomrule
\end{tabular}
\caption{Valori memorizzati da STERA 30M per componente. Tutti i valori sono memorizzati e valutati in FP32.}
\label{tab:componenti-stera30m}
\end{table}

\paragraph{Al crescere del modello.} La \autoref{tab:risparmio-crescita} applica la \autoref{prop:risparmio-modello} a modelli con la stessa forma di STERA 30M e dimensione crescente. Senza low-rank il risparmio supera il $78\%$ per ogni dimensione considerata. Con il rango fissato a 96 cresce con il modello: è del $63{,}3\%$ per $d=768$ e supera il $76\%$ per $d=4096$. Con il rango proporzionale alla dimensione resta quasi costante; scende leggermente perché il peso dell'embedding, che risparmia più dei layer, diminuisce, e tende al valore per layer del $61{,}9\%$.

\begin{table}[H]
\centering
\small
\begin{tabular}{rrrrr}
\toprule
$d$ & senza low-rank & $r=96$ & $r=d/8$ & layer con $r=d/8$ \\
\midrule
512 & $78{,}5\%$ & $56{,}2\%$ & $63{,}6\%$ & $61{,}6\%$ \\
768 & $78{,}9\%$ & $63{,}3\%$ & $63{,}3\%$ & $61{,}7\%$ \\
1024 & $79{,}1\%$ & $67{,}1\%$ & $63{,}1\%$ & $61{,}7\%$ \\
2048 & $79{,}7\%$ & $73{,}3\%$ & $62{,}6\%$ & $61{,}8\%$ \\
4096 & $80{,}1\%$ & $76{,}7\%$ & $62{,}2\%$ & $61{,}8\%$ \\
\bottomrule
\end{tabular}
\caption{Risparmio rispetto alla forma densa di modelli con sei layer, vocabolario di \num{50304} token e \num{2048} posizioni, al variare della dimensione $d$ e del rango. Le prime tre colonne si riferiscono al modello intero, l'ultima al solo layer. Per $d=768$ e $r=96$ si ritrova STERA 30M.}
\label{tab:risparmio-crescita}
\end{table}

\section{Riepilogo}

Una proiezione STERA è formata da quattro componenti. Le coordinate dei blocchi RowMask23 sono la parte strutturata: un sottospazio lineare di dimensione \num{3176} su \num{16384} per blocco, quindi una rappresentazione esatta e cinque volte più compatta, con matrici in generale a rango pieno. Scale e guadagni costano meno dello $0{,}5\%$ dei valori e regolano l'ampiezza di blocchi, righe e colonne; il termine low-rank accoppia i blocchi con poche direzioni senza struttura. Il bilancio complessivo separa ciò che il formato garantisce, l'$80{,}61\%$ di risparmio sulla parte strutturata, da ciò che dipende dalle scelte di progetto, cioè il costo $1{,}5\,\kappa$ del low-rank per layer.

Nel prossimo capitolo vediamo come una proiezione STERA esegue il forward e la backpropagation, cioè come si calcola e come si addestra.
